\subsection{Outer measure of an open set}

Given an open set \(G \subset X\) , its characteristic function \(\varphi_{G}\) is lower semicontinuous on X (GT, IV, §6, No.~2, Cor. of Prop. 1). We may therefore make the following definition:

DEFINITION 2. --- Given a positive measure \(\mu\) on X, for every open set \(G \subset X\) the upper integral \(\mu^{*}(\varphi_{G})\) is called the outer measure of G and is denoted \(\mu^{*}(G)\) .

The outer measure of an open set G is thus a number \(\geqslant0\) , finite or equal to \(+\infty\) . Clearly \(\mu^{*}(\varnothing)=0\) . Moreover, \(\mu^{*}(X)=\|\mu\|\) , as is shown by formula (22) of Ch. III, §1, No.~8.

PROPOSITION 5. --- The outer measure of a relatively compact open set G is finite.

For, there exists in this case a function \(f \in \mathcal{K}_+\) such that \(\varphi_{\mathrm{G}} \leqslant f\) (Ch. III, §1, No.~2, Lemma 1), whence

\[
\mu^ {*} (\mathrm{G}) = \mu^ {*} (\varphi_ {\mathrm{G}}) \leqslant \mu^ {*} (f) = \mu (f) <   + \infty .
\]

An open set of finite outer measure is not always relatively compact (Exer. 3).

PROPOSITION 6. --- If \(G_1\) and \(G_2\) are two open sets such that \(G_1 \subset G_2\) , then \(\mu^*(G_1) \leqslant \mu^*(G_2)\) .

For, the relation \(\mathbf{G}_1\subset \mathbf{G}_2\) is equivalent to \(\varphi_{\mathrm{G_1}}\leqslant \varphi_{\mathrm{G_2}}\)

PROPOSITION 7. --- Let G be a set of open subsets of X that is directed for the relation ⊂; then

\[
\mu^ {*} \left(\bigcup_ {G \in \mathfrak {G}} G\right) = \sup _ {G \in \mathfrak {G}} \mu^ {*} (G).\tag{5}
\]

The functions \(\varphi_{\mathbf{G}}\) form a directed set (for \(\leqslant\) ) in \(\mathcal{I}_{+}\) and their upper envelope is the characteristic function of the union of the sets \(\mathbf{G} \in \mathfrak{G}\) ; the proposition is thus a consequence of Th. 1.

PROPOSITION 8. --- Let \((G_{\iota})_{\iota \in I}\) be any family of open sets; then

\[
\mu^ {*} \left(\bigcup_ {\iota \in I} G _ {\iota}\right) \leqslant \sum_ {\iota \in I} \mu^ {*} (G _ {\iota}).\tag{6}
\]

Moreover, if the \(G_{\iota}\) are pairwise disjoint then

\[
\mu^ {*} \left(\bigcup_ {\iota \in I} G _ {\iota}\right) = \sum_ {\iota \in I} \mu^ {*} (G _ {\iota}).\tag{7}
\]

For, if \(G = \bigcup_{\iota \in I} G_{\iota}\) then \(\varphi_{G} = \sup_{\iota \in I} \varphi_{G_{\iota}} \leqslant \sum_{\iota \in I} \varphi_{G_{\iota}}\) ; when the \(G_{\iota}\) are pairwise disjoint, \(\varphi_{G} = \sum_{\iota \in I} \varphi_{G_{\iota}}\) ; the proposition is therefore a consequence of Props. 2 and 3.

Example. --- Let X = R and let \(\mu\) be Lebesgue measure on R (Ch. III, §1, No.~3); we are going to determine the outer measure of an open interval G = {]}a, b{[} ( \(-\infty \leqslant a < b \leqslant +\infty\) ). Suppose first that a and b are finite. For every function f in \(K_{+}\) such that \(f \leqslant \varphi_{G}\) , we have, by the theorem of the mean (FRV, II, §1, No.~5, Prop. 6),

\[
\int_ {- \infty} ^ {+ \infty} f (x) d x = \int_ {a} ^ {b} f (x) d x \leqslant b - a,
\]

whence \(\mu^{*}(\mathrm{G})\leqslant b - a\) . On the other hand, for every \(\varepsilon >0\) there exists a function \(f\in \mathcal{K}_{+}\) such that \(f\leqslant \varphi_{\mathrm{G}}\) and \(f(x) = 1\) for \(a + \varepsilon \leqslant x\leqslant b - \varepsilon\) whence \(\mu^{*}(G) \geqslant b - a - 2\varepsilon\) ; since \(\varepsilon\) is arbitrary, \(\mu^{*}(G) = b - a\) ; in other words, the outer measure of G is equal to its length. This result extends at once to the case that G is an unbounded open interval, since it then contains bounded open intervals of arbitrarily large length; thus \(\mu^{*}(G) = +\infty\) in this case.

Now let G be any open set in R; G is the union of a countable set (finite or infinite) of pairwise disjoint open intervals \({]}a_{k}, b_{k}{[}\) (GT, IV, §2, No.~5, Prop. 2), consequently

\[
\mu^ {*} (\mathrm{G}) = \sum_ {k} (b _ {k} - a _ {k})
\]

(Prop. 8); in other words:

PROPOSITION 9. --- For Lebesgue measure on R, the outer measure of an open set in R is equal to the sum of the lengths of its connected components.

Note in particular that if G is an open set in R such that \(\mu^{*}(G)=0\) , then G is empty.

